\chapter{数学的思考是什么}

\section{Thinking vs.\ Calculating}
思考不等于计算。计算是沿着既定规则一步步执行；思考则要从具体方法中看出背后的客观规律。前者回答"怎么做"，后者回答"为什么"。

这本书不打算讲高深的定理。我想以一个同样在数学入门阶段的数学爱好者的身份，用一些简单的例子，分享自己遇到问题时是怎么想的。如果这些想法能对同样在路上的读者有一点启发，就足够了。

我们可以拿一个简单的例子来说，求解方程：
\[
\frac{1}{x^2+x-2} + \frac{1}{x^2+7x+10} = 2
\]

如果只靠计算，通常做法是通分，把分子乘开，你会得到一个四次多项式，然后面对一个枯燥的四次方程。这一步不需要动脑，但这真的是你想要的数学吗？

停下来想一想，观察分母，他们可以分别分解为 \((x+2)(x-1)\) 和 \((x+5)(x+2)\) 即 \((x-1)(x+2)\) 和 \((x+5)(x+2)\)，对于根 \(1\) 和 \(-5\) 的中点恰好是 \(-2\)，那是否说明，这意味着这两个二次式关于 \(x=-2\) 对称呢？

于是，令 \(x = t - 2\)，将对称中心移到原点：
\[
\frac{1}{t^2-3t} + \frac{1}{t^2+3t} = 2
\]

此时的结构豁然开朗，通分后分子相消，方程就为 \(\frac{2}{t^2-9}=2\)，解得 \(t = \pm\sqrt{10}\)。

再将 \(t\) 翻回 \(x\)，有
\[
x = -2 + \sqrt{10} \quad \text{或} \quad x = -2 - \sqrt{10} 
\]。

留意一个前提：\(t \neq 0\) 也就是说如果 \(t = 0\)，即 \(x = -2\)，此时原方程的分母会变成 \(0\)，方程没有意义。因此 \(x = -2\) 是原方程的增根，需要舍去。

最后，方程的解就为 \(x = -2 \pm \sqrt{10}\)

求根公式和通分规则告诉你该怎么做，但是他们并没有告诉你为什么应该是这样子的。
刚刚使用的一个“另类方法”均值换元，是由一个建立在观察之上的对称性思考出来的，再有化简过程中对分母不为 \(0\) 的警惕性，就是我要说的思考。

\section{Thinking vs.\ Rote}
上一节当中，我们用均值换元的思想解释了一个静态多项式方程。但当变量开始发生变化时，我们需要一种处理动态函数变化的工具。让我们把目光投向微积分！

\[
\boxed{f'(x)=\lim_{\Delta x\to 0}\frac{f(x+\Delta x)-f(x)}{\Delta x}}
\]

这是一个很漂亮直白的式子，逻辑上干脆利落，很容易理解其本质，但对于具体的应用，如果对直觉上不太敏感，怕是难以知其应用，那我们不妨试验一下，丢几个基础的多项式进去，看看她能反馈出什么东西来。

带入一个经典的一次函数 \(f(x) = kx + b\)。代入定义，有 

\[
f'(x) = \lim_{\Delta x \to 0} \frac{\left[ k(x + \Delta x) + b \right] - (kx + b)}{\Delta x} = k
\] 

我们发现，结果是 \(k\)。且无论 \(b\) 值如何变化， \(f'(x)\) 值都永远是 \(k\)。


我们再对经典的一个二次函数 \(f_2(x) = ax^2 + bx + c\) 研究。代入定义，我们先单独计算分子部分的差值：

\[
f(x + \Delta x) - f(x) = [a(x + \Delta x)^2 + b(x + \Delta x) + c] - (ax^2 + bx + c)
\]

展开平方项并化简，有

\[
2ax \Delta x + a(\Delta x)^2 + b\Delta x
\]

现在，我们把分子项重新放回式子，就为

\[
f'(x)=\lim_{\Delta x\to 0}\frac{2ax \Delta x + a(\Delta x)^2 + b\Delta x}{\Delta x}
\]

上下同时约去一个 $\Delta x$，最终结果就是
\[
f'(x) = 2ax + b
\] 

观察到，二次函数的导数变成了一次函数，如果你令 $2ax + b = 0$ 就会发现你得到了一个公式 $x = -b/2a$，正是抛物线的对称轴！原来，对称轴的几何意义，就是斜率为 $0$ 的地方！

既然规律如此一致，那我们可以从直觉上得到一个结论（经验证是正确的，这里只是为了方便）。我们不用再一遍遍套极限式子去算了，我们总结出：

\[
(x^n)' = n x^{n-1}
\]

再回想下刚刚的知识，你会意识到一个致命的问题，逆向推导时的常数项是被回收的，也就是说逆向推导出来的原函数带有不确定性，我们可以增添一个常数项带入计算得到结果。同时，我们意识到，看透了事务变化的结构，并不意味着掌握了事情的全貌。